\chapter{Service Discovery: Eureka}
\label{ch:eureka}

\section{The problem}

The gateway must forward \code{/api/courses} to course-service. To which
address? In Compose it could hard-code \code{course-service:8082} ---
but then every new replica, every moved service, every environment is a
config change. Service discovery replaces addresses with \emph{names
resolved at runtime}: services register themselves in a registry; clients
ask the registry.

\section{How Eureka works}

Eureka (from Netflix, the ancestral Spring Cloud stack) is a REST service
holding a registry of instances:

\begin{itemize}
  \item \textbf{Registration.} At startup each client POSTs its identity:
  app name (\code{spring.application.name}), host, port, health URL.
  \item \textbf{Heartbeats.} Every 30\,s (default) the client renews its
  lease. Miss a few renewals and the instance is evicted.
  \item \textbf{Client-side caching.} Consumers fetch the registry and
  cache it (default refresh 30\,s), then load-balance \emph{locally}
  across the instances they know. The registry being briefly down does
  not stop traffic --- callers keep using their cache.
  \item \textbf{Self-preservation.} If too many heartbeats vanish at once,
  Eureka assumes a network problem rather than mass death and stops
  evicting. Fewer surprises during network blips; staler data during real
  outages.
\end{itemize}

The consumer side in this project is the gateway's \code{lb://course-service}
URI (Chapter~4): ``look this name up in the registry, pick an instance.''

\begin{figure}[h]
\centering
\begin{tikzpicture}[node distance=5mm and 12mm]
  \node[boxdark] (eureka) {Eureka\\\scriptsize discovery-server:8761};
  \node[box, below left=of eureka] (a) {course-service\\\scriptsize instance 1};
  \node[box, below=of eureka] (b) {course-service\\\scriptsize instance 2};
  \node[box, below right=of eureka] (gw) {api-gateway};
  \draw[flow] (a) -- node[lbl,left]{register + heartbeat} (eureka);
  \draw[flow] (b) -- (eureka);
  \draw[flowdash] (gw) -- node[lbl,right]{fetch registry} (eureka);
  \draw[flow] (gw) -- node[lbl,below]{\scriptsize lb:// picks one} (b);
\end{tikzpicture}
\caption{Registration, heartbeat, and client-side load balancing.}
\end{figure}

\section{The server, and the one k8s-critical setting}

The discovery-server module is Eureka with almost no configuration ---
a \code{@EnableEurekaServer} annotation and a port. The interesting
configuration is on every \emph{client}, and one line of it earns its own
section because it decides whether the whole system works on Kubernetes:

\begin{lstlisting}[language=yaml]
- name: EUREKA_INSTANCE_PREFER_IP_ADDRESS
  value: "true"
\end{lstlisting}

By default a client registers its \emph{hostname}. On your laptop that is
fine. Inside a Kubernetes pod, the hostname is the \emph{pod name} ---
\code{course-service-7d9f6c-x2vlq} --- which is not resolvable by other
pods. Every registration would be an address no one can dial.
Pod \emph{IPs}, by contrast, are routable across the whole cluster.
So every Deployment in \code{deploy/k8s/base/} sets this flag: register
the IP, not the name.

\begin{pitfall}[everything registers, nothing can call anything]
Symptom: the Eureka dashboard shows all services UP, but every
gateway-routed request fails with \code{UnknownHostException}. Cause: the
registered hostnames are pod names (see above). This is the single most
common Eureka-on-Kubernetes failure, and it is invisible until the first
cross-service call.
\end{pitfall}

\begin{decision}[why Eureka at all, when Kubernetes has Services]
Kubernetes already gives every Service a DNS name
(\code{course-service:8082}) with built-in load balancing --- discovery is
native to the platform. Eureka here is deliberate résumé-driven
redundancy: the classic stack is still widespread in industry, and
migrating \emph{off} it (Chapter~27: replace \code{lb://} with plain
Service DNS, delete two deployables) is itself the kind of modernization
work enterprises pay for. Building both stacks teaches the migration.
\end{decision}

\section{Outside the project}

\emph{Consul} (HashiCorp) is the main non-JVM alternative: agent-based,
language-agnostic, with health checks executed by the agent rather than
self-reported heartbeats, plus a key-value store that can replace the
config-server too. \emph{Kubernetes-native discovery} is the other pole:
no registry process at all --- the platform's controllers watch pod
readiness and update Service endpoints; ``registration'' is simply a pod
becoming Ready (Chapter~17 shows the machinery). Notice the pattern: the
more the platform provides, the less the application must carry.

\begin{hardware}
Eureka's conservative defaults (30\,s heartbeat, 30\,s client cache,
eviction after 90\,s) mean a crashed instance can receive traffic for up
to a minute. On one replica per service that hardly matters --- the caller
would fail anyway --- but it explains why rolling deployments in
Chapter~23 wait on Kubernetes readiness, not on Eureka status.
\end{hardware}

\begin{quiz}
\begin{enumerate}
  \item Trace what happens between course-service starting and the
  gateway successfully routing to it: which registrations, fetches, and
  caches are involved, and what is the worst-case delay?
  \item Why does \code{prefer-ip-address} matter on Kubernetes but not on
  Compose? (What does each platform make resolvable?)
  \item Self-preservation mode trades which property for which? Give a
  scenario where it helps and one where it hurts.
  \item If Eureka were removed and the gateway pointed at
  \code{http://course-service:8082} directly, what Kubernetes machinery
  replaces each Eureka feature?
\end{enumerate}
\end{quiz}
